\begin{lstlisting}[language=Bash, caption={Multipass cluster provisioning script for Scenario 1 (Vanilla Kubernetes)}, label={lst:multipass-scenario1}]
#!/usr/bin/env bash
# ==============================================================================
# Multipass Cluster Provisioning Script - Scenario 1 (Vanilla Kubernetes)
# ==============================================================================

echo "Deploying 1 Control Plane (Master) and 2 Worker Nodes via Multipass..."

# 1. Launch Virtual Machines with identical hardware constraints (2 vCPU, 4GB RAM)
multipass launch --name k8s-master --cpus 2 --memory 4G --disk 40G 22.04
multipass launch --name k8s-worker1 --cpus 2 --memory 4G --disk 30G 22.04
multipass launch --name k8s-worker2 --cpus 2 --memory 4G --disk 30G 22.04

# 2. Extract Master IP for clustering
MASTER_IP=$(multipass info k8s-master --format json | jq -r '.info["k8s-master"].ipv4[0]')
echo "Master node successfully provisioned at IP: $MASTER_IP"

# 3. Provision K3s Server on Master Node
multipass exec k8s-master -- bash -c "
  curl -sfL https://get.k3s.io | INSTALL_K3S_EXEC=\"--write-kubeconfig-mode 644 --node-ip=$MASTER_IP\" sh -
"

# 4. Extract cluster token from Master
NODE_TOKEN=$(multipass exec k8s-master -- sudo cat /var/lib/rancher/k3s/server/node-token)

# 5. Provision K3s Agent on Worker Nodes pointing to the Master IP
for NODE in k8s-worker1 k8s-worker2; do
  echo "Installing K3s Agent on $NODE..."
  WORKER_IP=$(multipass info $NODE --format json | jq -r ".info[\"$NODE\"].ipv4[0]")
  multipass exec $NODE -- bash -c "
    curl -sfL https://get.k3s.io | K3S_URL=https://$MASTER_IP:6443 K3S_TOKEN=$NODE_TOKEN INSTALL_K3S_EXEC=\"--node-ip=$WORKER_IP\" sh -
  "
done

echo "Vanilla Kubernetes cluster successfully created!"
\end{lstlisting}